

\documentclass[notes, 11pt]{beamer} % THIS
% \documentclass[notes, 11pt]{beamer}
% \documentclass[notes, 11pt, aspectratio=169]{beamer}

% \usepackage{helvet}
% \usepackage[default]{lato}
% \usepackage{array}


% \documentclass[handout]{beamer}
% \usepackage{pgfpages}
% \pgfpagesuselayout{4 on 1}[a4paper,border shrink=5mm]

\usepackage{bbm}
\usepackage{amsthm}
\usepackage{amssymb}
\usepackage{xcolor}
\usepackage{graphicx}
\usepackage{epstopdf}
\usepackage{booktabs}
\usepackage[percent]{overpic}

\usepackage{sgamevar, tikz} % Game theory packages
\usetikzlibrary{trees, calc} % For extensive form games
\usepackage{subfig} % Manipulation and reference of small or sub figures and tables

\setbeamertemplate{theorems}[numbered] % to number
% \usetheme{Malmoe}
\usetheme{default}
% \usecolortheme{whale}
\usecolortheme{rose}


\newtheorem{defi}{Definition}
\newtheorem{teo}{Theorem}
\newtheorem{remi}{Remark}
\newtheorem{exi}{Example}
\newtheorem{propi}{Proposition}
\newtheorem{proteo}{Proof}
\newtheorem{lemi}{Lemma}
\newtheorem{prolemi}{Proof}
\newtheorem{propropi}{Proof}
\newtheorem{coro}{Corollary}
\newtheorem{proci}{Algorithm}
\newtheorem{ass}{Assumption}

\setbeamertemplate{footline}[frame number]

\usepackage{caption}
\captionsetup[figure]{labelformat=empty}%

\newenvironment{wideitemize}{\itemize\addtolength{\itemsep}{10pt}}{\enditemize}
% \usepackage{enumitem}

\usepackage[spanish]{babel}

% \renewcommand{\baselinestretch}{1}
% \expandafter\def\expandafter\insertshorttitle\expandafter{%
%   \insertshorttitle\hfill%
%   \insertframenumber\,/\,\inserttotalframenumber}

\DeclareMathOperator*{\argmax}{arg\,max}

\begin{document}
\title{Microeconom\'ia Aplicada}   
\subtitle{Teor\'ia de Juegos: Juegos Est\'aticos con Informacion Incompleta}   
\author[Paz y Mino ]{Jos\'e Manuel Paz y Mi\~no}
\institute[shortinst]{FEN UCHILE}
\date{Oto\~no 2025} 

\frame{\titlepage} 

\begin{frame}{Introducci\'on}
	\begin{wideitemize}
		\item<1-> No siempre jugadores conocen toda la informacion de un juego
		\item<2-> Ejemplos
		\begin{wideitemize}
			\item<3->[--] Costos rivales.
			\item<4->[--] Valores rival de objeto en subasta.
			\item<5->[--] Demanda por un producto.
		\end{wideitemize}
		\item<6-> Como m\'aximo, jugadores tienen estimados sobre estos valores.

	\end{wideitemize}
\end{frame}


\begin{frame}{Elementos}
	\begin{wideitemize}
		\item<1-> N\'umero finito de jugadores $N$.
		\item<2-> Set finito de estados de la naturaleza $\Omega$ (con elemento $w$).
		\item<3-> Para cada jugador $i \in N$
		\begin{wideitemize}
			\item<4-> Set de acciones $A_{i}$
			\item<5-> Set de se\~nales (o set de tipos) que pueden ser observadas por cada jugador $T_{i}$ y una funci\'on de se\~nal $\tau_{i}:\Omega \rightarrow T_{i}$.
			\item<6-> Probabilidad $p_{i}$ en $\Omega$ (probabilidad prior) tal que $p_{i}(\tau_{i}^{-1}(t_{i}))$ para todo $t_{i} \in T_{i}$.
			\item<7-> Relaci\'on de preferencia $\succsim_{i}$ sobre $A \times \Omega$.
		\end{wideitemize}
	\end{wideitemize}
\end{frame}

\begin{frame}{Elementos}
	\begin{wideitemize}
		\item<1-> $\Omega$ recoge todas las posibilidades sobre las que se tiene informaci\'on incompleta.
		\item<2-> Note que puede haber una diferencia entre el estado $w$ y la se\~nal que se revela a cada jugador $t_{i}$.
		\item<3-> Tambi\'en jugadores pueden tener distintas probabilidades prior (pueden diferir en la probabilidad que piensan los estados de la naturaleza ocurren).
		\item<4-> Decimos que cuando el juego ocurre cada jugador conoce su tipo.
	\end{wideitemize}
\end{frame}

\begin{frame}{Elementos}
	\begin{wideitemize}
		\item<1-> No se pueden asumir los estados por separado (como juegos separados). En un estado especifico, hay que tener creencias sobre lo que harian otros jugadores en otros estados. 
		\item<2-> Entonces un equilibrio de Nash de un juego Bayesiano $<N, \Omega, (A_{i}), (T_{i}), (\tau_{i}), (p_{i}), (U_{i})>$ es un equilibrio de Nash del juego estrategico $G^{*}$ en donde para cada $i \in N$ y posible se\~nal $t_{i} \in T_{i}$, hay un jugador referido como $(i,t_{i})$ (\textit{tipo }$t_{i}$ \textit{ del jugador }$i$).
		\item<3-> Entonces una estrategia para el jugador $i$ es una funci\'on $s_{i}: T_{i} \rightarrow A_{i}$ y un perfil de estrategia $s = (s_{i})$
	\end{wideitemize}
\end{frame}

\begin{frame}{Equilibrio de Nash del juego bayesiano}
	\begin{wideitemize}
		\item<1-> Un perfil de estrategias $s^{*}$ es un equilibrio de Nash del juego bayesiano si para cada $i \in N$ y $t_{i} \in T_{i}$,
		\begin{align*}
			\tilde{u}(s^{*};t_{i}) \geq \tilde{u}_{i}(a_{i},s_{-i}^{*};t_{i})
		\end{align*} para todo


	\end{wideitemize}
\end{frame}


\begin{frame}
\maketitle
\end{frame}

\end{document}
